% 2de graad, Stelsels (reeks Vergelijkingen).
% De Hooch: een vrouw met een weegschaal; twee voorwaarden die tegelijk moeten kloppen.
% Motief: de rechten y = x + 1 en y = -x + 5 met hun snijpunt (2, 3), de oplossing van het stelsel.
\def\omslagbestand{omslag-goudweegster.jpg}
\def\omslagbron{Pieter de Hooch, \textit{De goudweegster}, ca. 1664 (publiek domein)}
\def\omslagrand{10.7 + 0.3*sin(\x*360/14)}
\def\omslagmotief{%
  \randlijn
  \begin{scope}[shift={(14.8,4.6)},scale=.8]
  \draw[white,line width=.7pt,-{Stealth[length=2mm]}] (-.3,0) -- (6,0);
  \draw[white,line width=.7pt,-{Stealth[length=2mm]}] (0,-.3) -- (0,6);
  \draw[white,line width=1.4pt] (0,1) -- (4.6,5.6);
  \draw[white,line width=1.4pt] (0,5) -- (5,0);
  \draw[white,line width=.6pt,dashed] (2,0) -- (2,3) -- (0,3);
  \fill[fel,draw=accent,line width=.7pt] (2,3) circle (2.2mm);
  \node[text=fel,anchor=west] at (2.25,3.2) {$(2,3)$};
  \end{scope}
}
